\section{Matching to short-distance amplitudes}
\label{sec:matching}

The matching map from the imported kernel and the sector data to the rare
kaon observables factorizes through the exact TFPT CKM matrix:
\begin{equation}
\label{eq:matching-map}
\mathcal M_{\mathrm{RK}}
:
\bigl(\text{TFPT kernel},\;\mathcal D_{\mathrm{RK}}\bigr)
\;\xrightarrow{\;\text{transport}\;}\;
(V_{\mathrm{CKM}})
\;\xrightarrow{\;\text{SD matching}\;}\;
\bigl(\BR(K^+),\,\BR(K_L)\bigr).
\end{equation}

\begin{definition}[Standard short-distance branching ratios]
\label{def:br-formulas}
We use the standard modern short-distance formulas \cite{BrodGorbahnStamou}:
\begin{equation}
\label{eq:br-kplus}
\BR(K^+\!\to\pi^+\nu\bar\nu)
=
\kappa_+\,(1+\Delta_{\mathrm{EM}})
\left[
\left(\frac{\im\lambda_t}{\lambda^5}\,X_t\right)^{\!2}
+
\left(
\frac{\re\lambda_c}{\lambda}\,(P_c+\delta P_{c,u})
+
\frac{\re\lambda_t}{\lambda^5}\,X_t
\right)^{\!2}
\right],
\end{equation}
\begin{equation}
\label{eq:br-kl}
\BR(K_L\to\pi^0\nu\bar\nu)
=
\kappa_L\,r_{\epsilon_K}
\left(\frac{\im\lambda_t}{\lambda^5}\,X_t\right)^{\!2},
\end{equation}
where $\lambda_t:=V_{ts}^*V_{td}$, $\lambda_c:=V_{cs}^*V_{cd}$,
$\lambda:=|V_{us}|=s_{12}$,
$X_t$ is the top-quark Inami--Lim function including NLO QCD and
electroweak corrections, $P_c$ is the charm contribution at NNLO QCD
and NLO EW, $\delta P_{c,u}$ is the long-distance charm--up
interference correction, and $r_{\epsilon_K}$ accounts for indirect
CP violation in the $K_L$ state.
\end{definition}

\begin{tcolorbox}[colback=gray!4, colframe=gray!60, breakable,
    title={\textbf{Short-distance input parameters}}, fonttitle=\bfseries]
\begin{itemize}[leftmargin=1.6em, itemsep=0.3em]
\item $X_t = 1.462\pm 0.017$ \quad (NLO QCD + NLO EW, at
      $m_t^{\mathrm{pole}}=\SI{172.6}{GeV}$, $M_W=\SI{80.36}{GeV}$);
\item $P_c = (0.2255/\lambda)^4\,(0.3604\pm 0.0087)$ \quad (NNLO QCD + NLO EW);
\item $\delta P_{c,u}=0.04\pm 0.02$ \quad (long-distance charm--up interference);
\item $\Delta_{\mathrm{EM}}=-0.003$ \quad (electromagnetic correction to $K^+$);
\item $\kappa_+=\num{5.173e-11}\,(\lambda/0.225)^8$,\;
      $\kappa_L=\num{2.231e-10}\,(\lambda/0.225)^8$;
\item $|\epsilon_K|=(\num{2.228}\pm\num{0.011})\times 10^{-3}$.
\end{itemize}
\end{tcolorbox}

At the present level of precision, the small indirect-CP-violation
correction to the neutral mode is included through \cite{BrodGorbahnStamou}
\begin{equation}
\label{eq:r-epsilon-K}
r_{\epsilon_K}
=
1
-\frac{\sqrt{2}\,|\epsilon_K|}
{1+P_c/(A^2 X_t)-\rho/\eta},
\qquad
\rho=\frac{\bar\rho}{1-\lambda^2/2},
\quad
\eta=\frac{\bar\eta}{1-\lambda^2/2}.
\end{equation}
On the TFPT branch, with $\bar\rho=0.1374$, $\bar\eta=0.3509$,
$A=0.8624$, $\lambda=0.22438$, this yields
\begin{equation}
\label{eq:r-epsilon-K-value}
r_{\epsilon_K}^{\mathrm{TFPT}}\approx 0.997.
\end{equation}

\begin{tcolorbox}[colback=blue!3, colframe=blue!50!black, breakable,
    title={\textbf{TFPT-fixed rare kaon prediction}}, fonttitle=\bfseries]
\begin{proposition}[TFPT-fixed CKM branching ratios with standard SD input]
\label{prop:br-prediction}
Insert the exact TFPT CKM products $\lambda_t$, $\lambda_c$ from
\Cref{eq:lambda-t-exact,eq:lambda-c-exact} into
\Cref{eq:br-kplus,eq:br-kl} with the short-distance parameters above.
TFPT fixes the CKM point without a flavor fit.
The short-distance functions $X_t$, $P_c$, $\delta P_{c,u}$, the
electromagnetic correction $\Delta_{\mathrm{EM}}$, and the hadronic
normalization constants $\kappa_\pm$ are standard external inputs.
The factor $r_{\epsilon_K}$ is derived from $|\epsilon_K|$ and the
TFPT CKM point via \Cref{eq:r-epsilon-K}.
The numerical evaluation yields
\begin{equation}
\label{eq:br-values}
\boxed{
\BR(K^+\!\to\pi^+\nu\bar\nu)
=
\num{9.40e-11},
\qquad
\BR(K_L\to\pi^0\nu\bar\nu)
=
\num{3.47e-11}.
}
\end{equation}
\end{proposition}
\end{tcolorbox}

\begin{proof}[Evaluation chain]
From \Cref{eq:lambda-t-exact,eq:lambda-c-exact} with
$\lambda=s_{12}=\num{0.22438}$:
\[
\frac{\im\lambda_t}{\lambda^5}
=
\frac{\num{1.522e-4}}{(0.22438)^5}
=
\num{0.2856},
\qquad
\frac{\re\lambda_t}{\lambda^5}
=
\frac{-\num{3.558e-4}}{(0.22438)^5}
=
\num{-0.6677},
\]
\[
\frac{\re\lambda_c}{\lambda}
=
\frac{-\num{0.2183}}{0.22438}
=
\num{-0.9729}.
\]
With $X_t=1.462$, $P_c=(0.2255/0.22438)^4\times 0.3604=\num{0.3625}$,
$\delta P_{c,u}=0.04$, $\Delta_{\mathrm{EM}}=-0.003$,
$\kappa_+=\num{5.173e-11}\times(0.22438/0.225)^8=\num{4.98e-11}$,
$\kappa_L=\num{2.231e-10}\times(0.22438/0.225)^8=\num{2.147e-10}$,
$r_{\epsilon_K}^{\mathrm{TFPT}}=0.997$:

\medskip
\noindent\textbf{Charged mode:}
\[
\BR(K^+)
=
\kappa_+'\times\bigl[
(\num{0.4175})^2
+
(\num{-1.3682})^2
\bigr]
=
\num{4.97e-11}\times\num{1.891}
=
\num{9.40e-11},
\]
where $\kappa_+'=\kappa_+(1+\Delta_{\mathrm{EM}})=\num{4.97e-11}$,
the CP-violating piece is $\im\lambda_t\,X_t/\lambda^5=\num{0.4175}$,
and the CP-conserving piece is
\[
\frac{\re\lambda_c}{\lambda}(P_c+\delta P_{c,u})
+\frac{\re\lambda_t}{\lambda^5}X_t
=
\num{-0.3917}-\num{0.9765}
=
\num{-1.3682}.
\]

\noindent\textbf{Neutral mode:}
\[
\BR(K_L)
=
\num{2.147e-10}\times 0.997
\times(0.2856\times 1.462)^2
=
\num{3.47e-11}.
\]
\end{proof}

\begin{remark}
The $K_L$ channel is purely CP violating and therefore proportional to
$(\im\lambda_t)^2$.
Since $\im\lambda_t\propto\bar\eta\propto\sin\delta_{\mathrm{CKM}}$,
the neutral-channel prediction is a direct test of the TFPT holonomy
phase $\delta_{\mathrm{CKM}}=\pi/3+3\lambda_C^2+O(\lambda_C^4)$.
The charged channel receives a sizeable charm admixture through
$P_c+\delta P_{c,u}$ and is therefore a less clean probe of the
CP phase alone, but a stronger test of the overall CKM norm.
\end{remark}


\subsection{Grossman--Nir correlation}

The model-independent Grossman--Nir (GN) bound provides a
constraint between the two channels that holds in any model
respecting isospin:
\begin{equation}
\label{eq:gn-bound}
\BR(K_L\to\pi^0\nu\bar\nu)
\;\le\;
4.3\;\BR(K^+\!\to\pi^+\nu\bar\nu).
\end{equation}
On the TFPT branch the ratio is
\begin{equation}
\label{eq:gn-ratio}
\frac{\BR(K_L)}{\BR(K^+)}
=
\frac{\num{3.47e-11}}{\num{9.40e-11}}
=
0.369,
\end{equation}
safely within the isospin corridor $(\le 4.3)$.

\begin{figure}[t]
\centering
\begin{tikzpicture}
\begin{axis}[
    width=0.88\linewidth,
    height=0.62\linewidth,
    xlabel={$\BR(K^+\!\to\pi^+\nu\bar\nu)\;\;[\times 10^{-11}]$},
    ylabel={$\BR(K_L\to\pi^0\nu\bar\nu)\;\;[\times 10^{-11}]$},
    xmin=0, xmax=22,
    ymin=0, ymax=55,
    grid=major,
    grid style={gray!20},
    axis lines=left,
    tick label style={font=\small},
    label style={font=\small},
    legend style={font=\scriptsize, at={(0.03,0.97)}, anchor=north west,
        draw=gray!40, fill=white, fill opacity=0.9},
    clip=false,
]

\addplot[domain=0:12.8, samples=80, thick, red!60!black, dashed,
    name path=gn] {4.3*x};
\addlegendentry{Grossman--Nir bound}

\addplot[only marks, mark=*, mark size=3.5pt, blue!70!black, fill=blue!70!black]
    coordinates {(9.40, 3.47)};
\addlegendentry{TFPT prediction}

\addplot[only marks, mark=diamond*, mark size=4pt, orange!80!black,
    fill=orange!60!white]
    coordinates {(8.60, 2.94)};
\addlegendentry{SM central (CKMfitter)}

\addplot[only marks, mark=square*, mark size=3pt, purple!70!black,
    fill=purple!30!white]
    coordinates {(13.0, 0)};
\addlegendentry{NA62 observation}
\draw[purple!50!black, thick] (axis cs:10.0,0.8) -- (axis cs:16.3,0.8);
\node[font=\scriptsize, purple!50!black] at (axis cs:13.0,2.2) {$13.0^{+3.3}_{-3.0}$};

\addplot[name path=killL, draw=none] coordinates {(7,0) (7,55)};
\addplot[name path=killR, draw=none] coordinates {(12,0) (12,55)};
\addplot[green!30, opacity=0.35] fill between[of=killL and killR];

\draw[green!50!black, thick, densely dotted] (axis cs:7,0) -- (axis cs:7,55);
\draw[green!50!black, thick, densely dotted] (axis cs:12,0) -- (axis cs:12,55);

\node[font=\scriptsize, green!40!black, rotate=90, anchor=south]
    at (axis cs:7.4,42) {test corridor};
\node[font=\scriptsize, green!40!black, rotate=90, anchor=south]
    at (axis cs:12.4,42) {test corridor};

\node[font=\scriptsize, blue!70!black, anchor=south west]
    at (axis cs:9.9,4.0) {$(9.40,\;3.47)$};

\fill[red!8, opacity=0.25] (axis cs:0,0) -- (axis cs:12.8,55.04) --
    (axis cs:22,55.04) -- (axis cs:22,0) -- cycle;
\node[font=\scriptsize, red!50!black, rotate=32, anchor=south]
    at (axis cs:6.0,27) {GN excluded};

\end{axis}
\end{tikzpicture}
\caption{Grossman--Nir plane for $K\to\pi\nu\bar\nu$. The dashed red line is
the model-independent isospin bound $\BR(K_L)\le4.3\,\BR(K^+)$; the region
above it is excluded. The blue dot marks the TFPT prediction; the orange
diamond shows the SM central value from CKMfitter. The purple square shows
the NA62 combined measurement $(13.0^{+3.3}_{-3.0})\times10^{-11}$
(charged channel only; the $K_L$ measurement has not yet reached SM sensitivity).
The green corridor $[7,12]\times10^{-11}$ is the operational TFPT test window
for the charged channel; it is not a formal confidence interval but an
approximate consistency band around the central prediction.}
\label{fig:gn-plane}
\end{figure}

Any future measurement placing the experimental ratio above $4.3$
would signal physics outside the SM operator basis and would
simultaneously invalidate the TFPT matching ansatz for this sector.
The predicted point in the $(K^+,K_L)$ plane is shown as a fixed target
in \Cref{fig:gn-plane}.


\section{Benchmark ledger}
\label{sec:benchmarks}

\begin{table}[ht]
\centering
\footnotesize
\begin{tabularx}{\linewidth}{@{}>{\raggedright\arraybackslash}p{0.17\linewidth}>{\raggedright\arraybackslash}Y>{\raggedright\arraybackslash}p{0.16\linewidth}>{\raggedright\arraybackslash}p{0.17\linewidth}>{\centering\arraybackslash}p{0.11\linewidth}@{}}
\toprule
Observable & TFPT value & Imported input & Sector input & Status \\
\midrule
$\BR(K^+\!\to\pi^+\nu\bar\nu)$
  & $\num{9.40e-11}$
  & $\phiz,\lambda_C,\delta_{\mathrm{CKM}}$
  & $X_t$, $P_c$, $\delta P_{c,u}$, $\Delta_{\mathrm{EM}}$, $\kappa_+$
  & prediction \\[6pt]
$\BR(K_L\to\pi^0\nu\bar\nu)$
  & $\num{3.47e-11}$
  & $\phiz,\lambda_C,\delta_{\mathrm{CKM}}$
  & $X_t$, $\kappa_L$, $|\epsilon_K|$
  & prediction \\[6pt]
$\BR(K_L)/\BR(K^+)$
  & $0.369$
  & CKM point
  & isospin ($\kappa_L/\kappa_+$)
  & GN test \\[6pt]
$\delta_{\mathrm{CKM}}$
  & $\num{1.198}$ rad
  & holonomy phase
  & ---
  & anchor \\[6pt]
$\lambda_C$
  & $\num{0.22438}$
  & $\sqrt{\phiz(1-\phiz)}$
  & ---
  & anchor \\
\bottomrule
\end{tabularx}
\caption{Benchmark ledger for the rare kaon note. The first two rows
are the primary predictions. The third is a derived GN-plane test.
The last two rows record the kernel anchors whose values propagate
into the branching ratios.}
\label{tab:benchmark-ledger}
\end{table}

\begin{table}[ht]
\centering
\small
\begin{tabularx}{\linewidth}{@{}>{\raggedright\arraybackslash}p{0.14\linewidth}>{\centering\arraybackslash}p{0.12\linewidth}>{\centering\arraybackslash}p{0.14\linewidth}>{\raggedright\arraybackslash}Y>{\centering\arraybackslash}p{0.09\linewidth}@{}}
\toprule
Parameter & TFPT & PDG\,'25 & Derivation chain & Pull \\
\midrule
$\lambda$ & $\num{0.22438}$ & $0.22501(68)$ & $\sqrt{\phiz(1-\phiz)}$ & $0.9\sigma$ \\
$A$ & $\num{0.8624}$ & $0.826^{+0.016}_{-0.015}$ & $\phiz/[\lambda^2(1+\lambda)]$ & $2.3\sigma$ \\
$\bar\rho$ & $\num{0.1374}$ & $0.1591(94)$ & exact $-V_{ud}V_{ub}^*/(V_{cd}V_{cb}^*)$ & $2.3\sigma$ \\
$\bar\eta$ & $\num{0.3509}$ & $0.3523(73)$ & exact $-V_{ud}V_{ub}^*/(V_{cd}V_{cb}^*)$ & $0.2\sigma$ \\
$\delta$ & $\num{1.198}$ rad & $1.147(26)$ rad & holonomy phase & $2.0\sigma$ \\
\bottomrule
\end{tabularx}
\caption{Exact TFPT CKM point versus the PDG\,2025 \cite{PDG2025} global fit.
The pulls for $A$, $\bar\rho$, and $\delta$ are at the $2\sigma$ level,
indicating genuine tension that the rare kaon measurements can resolve.}
\label{tab:wolfenstein}
\end{table}

The two branching-ratio rows depend on the imported CKM point
(from the kernel seed and the transport grammar) and on standard
electroweak loop functions (sector-specific, but not TFPT specific).
Every TFPT-specific content enters through $\lambda_C$,
$\delta_{\mathrm{CKM}}$, and the exact CKM matrix built from them.
The factor $r_{\epsilon_K}$ is not an independent external constant but
is derived from $|\epsilon_K|$ and the TFPT CKM point.


